\section{Scope, source audit, and exact status}

\subsection{What the two source pages actually assert}

The Wikipedia page gives one definition, three contour types, convergence
conditions, a differential equation, two residue expansions, conversion to and
from generalized hypergeometric functions, a polynomial repair formula,
parameter identities, derivatives, recurrences, scaling expansions, integral
formulas, transform kernels, and a long table of special cases.  MathWorld gives
the two sign conventions for the Mellin--Barnes integral, the four-list Wolfram
Language convention, a generalized fourth parameter $r$, and five special
values.  We treat a displayed equation, its domain clause, and its convergence
clause as separate claims.  This is important: several displayed identities
are correct only after adding hypotheses that the prose suppresses.

The machine-readable ledger \code{claims/coverage.csv} contains the complete
itemization.  Its status codes are:
\begin{center}
\begin{tabularx}{0.94\textwidth}{@{}cX@{}}
\toprule
\statusP & proved by a finite algebraic, coefficient, residue, or tail argument developed in this report;\\
\statusQ & proved under an explicit quantitative admissibility, Fubini, or transform certificate;\\
\statusC & source statement corrected or restricted (branch, boundary, resonance, or convergence);\\
\statusR & reduced to a reusable certificate theorem, but a broad ``all suitable functions'' assertion is not made;\\
\statusX & literal source wording is undefined or cannot be true on the advertised whole domain.\\
\bottomrule
\end{tabularx}
\end{center}
A status describes the human mathematical argument, not a completed Lean
object proof.

\subsection{Main constructive theorem}

Fix integers $0\le m\le q$ and $0\le n\le p$, finite parameter lists
$\mathbf a=(a_1,\ldots,a_p)$ and $\mathbf b=(b_1,\ldots,b_q)$, a branch of
$\Log$, and an argument $z\ne0$.  Define
\begin{equation}\label{eq:phi-minus}
 \Phi(s)=
 \frac{\prod_{j=1}^{m}\Gamma(b_j-s)
       \prod_{j=1}^{n}\Gamma(1-a_j+s)}
      {\prod_{j=m+1}^{q}\Gamma(1-b_j+s)
       \prod_{j=n+1}^{p}\Gamma(a_j-s)}.
\end{equation}
The traditional notation is
\begin{equation}\label{eq:main-definition}
 G_{p,q}^{m,n}\!\left(\left.\begin{matrix}\mathbf a\\\mathbf b\end{matrix}\right|z\right)
 =\frac1{2\pi\ii}\int_L \Phi(s)z^s\dd s.
\end{equation}
In iPRA, \eqref{eq:main-definition} is not a primitive definition.  A datum
must also supply a rational polygonal or vertical contour, pole-clearance and
branch certificates, and primitive-recursive error bounds.

\begin{theorem}[Certified equivalence of presentations]\label{thm:master}
For every fixed coded parameter datum satisfying one of the admissibility
contracts in \cref{sec:definitions,sec:convergence}, and every requested
precision $k$, there is a primitive-recursive procedure producing a rational
complex ball of radius at most $2^{-k}$ which contains each of the following
quantities whenever that presentation is admissible:
\begin{enumerate}[label=(\roman*)]
\item the truncated Mellin--Barnes integral with certified tails;
\item the right- or left-residue normal form, including confluent Laurent jets;
\item the local solution of the Meijer differential equation selected by the
      same Frobenius data.
\end{enumerate}
Whenever two presentations are simultaneously admissible on the same branch
domain, their two balls overlap.  In the nonconfluent right-loop regime the
normal form is the finite sum of generalized hypergeometric series in
\cref{thm:slater-right}; in the left-loop regime it is
\cref{thm:slater-left}.  All algorithms and all error schedules are primitive
recursive.  The proof uses only finite sums, products, bounded searches,
quantifier-free induction, and previously specified finite Gamma and residue
contracts.
\end{theorem}

The proof occupies \cref{sec:hypergeom,sec:definitions,sec:convergence}.  Its
value is logical as well as computational: all later identities can be proved
on whichever presentation makes the claim finite.  Parameter identities are
proved at the integrand level; local analytic identities are proved
coefficientwise; integral transforms are proved by finite Fubini estimates.
No unrestricted identity theorem is needed.

\subsection{Three levels that must not be conflated}

The accompanying files establish three different things.
\begin{enumerate}
\item The present report gives a human-readable constructive proof with explicit
      finite contracts and error reductions.
\item The Python programs perform exact-rational recurrence checks and
      high-precision numerical regression tests.  They are evidence against
      transcription errors, not proofs of universal statements.
\item The Lean file gives executable finite definitions and target theorem
      signatures.  It intentionally contains clearly marked placeholders for
      the analytic contracts and does not masquerade as an object-iPRA proof.
\end{enumerate}
The final ProveIt endpoint is a term of
\[
 []\;\vdash[\mathrm{PRAAx}]\;\varphi.
\]
A theorem about Lean's host type \code{Complex} or \code{Rat} is not that
endpoint.  Once such a derivation is built, ProveIt's checked theorem
\code{pra\_iff\_native} transfers it to the concise native scalar calculus.

\subsection{Corrections to literal web-page wording}

\begin{audit}[Fundamental systems]
The displayed $p\le q$ and $p\ge q$ lists are fundamental systems only under a
nonresonance hypothesis, for example $b_i-b_j\notin\Z$ in the first list and
$a_i-a_j\notin\Z$ in the second, together with nonzero leading coefficients.
At resonance, logarithmic Frobenius solutions or parameter derivatives replace
some listed functions.  The differential equation remains valid.
\end{audit}

\begin{audit}[``For all $x$'']
The formulas for $\log(1+x)$, inverse trigonometric functions, powers, and
Lerch transcendent require a branch.  The Heaviside formulas are pointwise
statements on the positive real axis away from $x=1$ unless a convention at the
jump is supplied.  The Lerch series has a pole or divergence at many advertised
points.  We prove branch-qualified or finite-observation versions rather than
the literal universal quantifier.
\end{audit}

\begin{audit}[Transform inversions]
The pages use ``suitable function'' and omit substantial convergence
conditions.  In weak arithmetic this cannot be left implicit.  We prove a
finite transform theorem from named endpoint, decay, quadrature, and Fubini
certificates.  Narain, Wimp, and the generalized Laplace formulas then become
instances.  We do not claim an unspecified maximal function-space theorem.
\end{audit}
